
We investigated whether retrieval quality diverges from pretraining quality. Our proof-of-concept experiments suggest ''yes'' directionally, while revealing how little we understand about what makes text retrievable.

This work makes three contributions. First, we provide an initial controlled empirical test suggesting retrieval-quality signals may diverge from pretraining fluency: corpus filtering trained on BEIR success examples achieved +10.6\% relative Recall@10 improvement over perplexity-based filtering in proof-of-concept validation. Second, we falsify the entity density hypothesis---demonstrating that NER-based factual density does not correlate with BEIR retrieval quality (ratio 0.973 < 1.15 threshold). Third, we document methodological challenges including corpus sampling issues affecting query coverage and the distinction between proof-of-concept validation and full-scale confirmation.

Retrieval-quality filtering shows initial promise in controlled experiments---classifiers can learn from BEIR annotations to select documents that may improve downstream retrieval performance at proof-of-concept scale. However, the quality signals they learn are not the entity-based factual density we hypothesized. The classifier learned something that appears to work in our controlled setting, but identifying that ''something'' remains an open question.

Our findings redirect research toward alternative mechanisms: query-document semantic alignment, answer-bearing sentence structures, conceptual density, or lexical diversity. The field needs theoretical frameworks for retrieval quality independent of pretraining paradigms. As RAG systems become ubiquitous, understanding what makes a corpus good for retrieval, independent of what makes it good for pretraining, will be essential.
